\chapter{Hermitian Observables, Spin, and the Probability Rule}

\lectureoverview{Quantum mechanics is a calculus for probabilities. To use it without replacing understanding by slogans, we need a compact mathematical chain: complex phases, normalized states, Hermitian observables, eigenvalues and eigenvectors, orthogonality, and the Born rule. The three Pauli matrices then turn that chain into concrete predictions for electron spin.}

\section{Complex Numbers in Polar Form}

Begin by removing one small source of algebraic friction. A complex number
\begin{equation}
  z=x+iy
\end{equation}
is a point in the complex plane, with \(x\) on the real axis and \(y\) on the imaginary axis. If the point has polar coordinates \((r,\theta)\), then
\begin{equation}
  x=r\cos\theta,
  \qquad
  y=r\sin\theta,
\end{equation}
and therefore
\begin{equation}
  z=r(\cos\theta+i\sin\theta)=re^{i\theta}.
  \label{eq:l03-polar-form}
\end{equation}
Euler's notation is useful because the angle-addition formulas become the exponential law
\begin{equation}
  e^{i\theta}e^{i\phi}
  =e^{i(\theta+\phi)}.
  \label{eq:l03-phase-product}
\end{equation}
Indeed, multiplying the two trigonometric forms gives
\begin{align}
& (\cos\theta+i\sin\theta)(\cos\phi+i\sin\phi)
\notag\\
&\qquad=\cos(\theta+\phi)+i\sin(\theta+\phi).
\end{align}
The familiar sum formulas are all contained in this one multiplication.

Complex conjugation reflects the point across the real axis:
\begin{equation}
  z^*=x-iy,
  \qquad
  (e^{i\theta})^*=e^{-i\theta}.
\end{equation}
Consequently,
\begin{equation}
  e^{i\theta}e^{-i\theta}=1.
\end{equation}
Numbers of the form \(e^{i\theta}\) lie on the unit circle and are called \emph{pure phases}. Every nonzero complex number is a positive modulus times a pure phase:
\begin{equation}
  z=|z|e^{i\theta},
  \qquad
  |z|^2=z^*z=x^2+y^2.
  \label{eq:l03-modulus-phase}
\end{equation}

\begin{classroomqa}
\lecturetimestamp{00:06:36}
\audiencequestion{Does changing \(\theta\) to \(-\theta\) correspond to reversing the direction of rotation in the complex plane?}
\lecturerresponse{Yes. Positive and negative phase angles trace the unit circle in opposite senses. Complex conjugation changes \(e^{i\theta}\) into \(e^{-i\theta}\), which is the reflected point and the inverse phase.}
\end{classroomqa}

That is the whole complex-number runway needed today: modulus, phase, conjugation, and the unit circle.

\section{Quantum Mechanics as a Probability Calculus}

We now turn to the postulates. A physical system is prepared in a \emph{state}; a measurement asks for the value of an \emph{observable}; quantum mechanics supplies probabilities for the possible answers.

States are vectors in a complex vector space. For a two-level spin system,
\begin{equation}
  \ket\psi=\alpha\ket{u}+\beta\ket{d},
  \qquad
  \ket{u}=\begin{pmatrix}1\\0\end{pmatrix},
  \qquad
  \ket{d}=\begin{pmatrix}0\\1\end{pmatrix}.
  \label{eq:l03-two-level-state}
\end{equation}
The coefficients determine the up/down probabilities,
\begin{equation}
  P(u)=|\alpha|^2,
  \qquad
  P(d)=|\beta|^2,
\end{equation}
so the state must be normalized:
\begin{equation}
  \braket{\psi}{\psi}
  =|\alpha|^2+|\beta|^2
  =1.
  \label{eq:l03-state-normalization}
\end{equation}
The dimension refers to the number of independent alternatives in the state space, not to the dimension of ordinary space. A system with three mutually exclusive basis states would use three orthonormal columns and three amplitudes.

Probabilities do not determine the amplitudes completely. Multiplication by a pure phase leaves a magnitude unchanged:
\begin{equation}
  |e^{i\theta}\alpha|^2=|\alpha|^2.
\end{equation}
The amplitudes therefore contain information beyond the probabilities for this one basis. Its role will become visible when different measurements and, later, time evolution are considered.

The basis states are orthonormal:
\begin{equation}
  \braket{u}{u}=\braket{d}{d}=1,
  \qquad
  \braket{u}{d}=\braket{d}{u}=0.
\end{equation}
Their orthogonality belongs to the complex space of states. It must not yet be confused with two arrows at right angles in ordinary three-dimensional space; that distinction will become decisive near the end of the lecture.

\section{Why Observables Are Hermitian}

Experimental results can be recorded as real numbers. We may package two real numbers into one complex number for convenience, but the pointer readings themselves can always be separated into real quantities. In finite dimensions, quantum observables are represented by a special class of linear operators: Hermitian matrices.

To define them, first transpose a matrix by interchanging rows and columns:
\begin{equation}
  (M^T)_{ij}=M_{ji}.
  \label{eq:l03-transpose}
\end{equation}
A real matrix satisfying \(M^T=M\) is symmetric. For complex matrices the corresponding operation is the Hermitian conjugate,
\begin{equation}
  M^\dagger=(M^T)^*,
  \qquad
  (M^\dagger)_{ij}=M_{ji}^*.
  \label{eq:l03-dagger}
\end{equation}

\begin{figure}[htbp]
\centering
\includegraphics[width=0.88\linewidth]{lecture_03_figure_02.png}
\caption{Transpose, symmetry, and the transition to Hermitian conjugation share one blackboard layout.}
\lectureframe{00:23:06}
\end{figure}

\begin{definition}
A matrix \(M\) is \emph{Hermitian} if
\begin{equation}
  M^\dagger=M,
  \qquad\text{equivalently}\qquad
  M_{ij}=M_{ji}^*.
  \label{eq:l03-hermitian}
\end{equation}
\end{definition}

The diagonal entries are real, since \(M_{ii}=M_{ii}^*\). Reflected off-diagonal entries are conjugate pairs. Thus the general two-by-two form is
\begin{equation}
  M=
  \begin{pmatrix}
    a&z\\
    z^*&b
  \end{pmatrix},
  \qquad a,b\in\mathbb R.
  \label{eq:l03-hermitian-two}
\end{equation}
Hermiticity is the matrix analogue of a reality condition; it does not require every entry separately to be real.

The first payoff follows from two elementary conjugation identities:
\begin{equation}
  \braket{B}{A}=\braket{A}{B}^*,
  \label{eq:l03-inner-conjugation}
\end{equation}
and, for any \(M\),
\begin{equation}
  \bra B M\ket A
  =
  \bigl(\bra A M^\dagger\ket B\bigr)^*.
  \label{eq:l03-matrix-element-conjugation}
\end{equation}
The expression \(\bra B M\ket A\) is a number, the \(BA\) matrix element of \(M\). If \(M\) is Hermitian and \(B=A\), then
\begin{equation}
  \bra A M\ket A
  =
  \bigl(\bra A M\ket A\bigr)^*.
\end{equation}
It is therefore real.

\begin{proposition}
For every Hermitian \(M\), the expectation value
\begin{equation}
  \expect{M}_A=\bra A M\ket A
  \label{eq:l03-expectation}
\end{equation}
is real in every normalized state \(\ket A\).
\end{proposition}

The probabilistic meaning of expectation value will be developed further, but its reality already explains why Hermitian operators are natural candidates for observables.

\begin{classroomqa}
\lecturetimestamp{00:33:22}
\audiencequestion{Does the angle-bracket notation for an average come from the bra--ket notation?}
\lecturerresponse{Yes. Dirac's bracket notation supplies the same brackets. In \(\bra A M\ket A\), the operator is sandwiched between the bra and ket representing the state, and the result is the expectation value.}
\end{classroomqa}

\begin{classroomqa}
\lecturetimestamp{00:34:31}
\audiencequestion{Can a physical quantity itself have an imaginary measured value? If measurements are real, why are complex numbers needed in quantum mechanics?}
\lecturerresponse{A set of measured results can always be recorded as real numbers, even if we choose to package a pair as one complex number. The sharper question is why the theory is driven to complex amplitudes. Hold that question: the answer is tied to reversible time evolution, which has not yet been introduced. Restricting the theory to real vectors would lose important possibilities.}
\end{classroomqa}

The mathematics is not an ornament. Without it, entanglement can be made to sound mysterious while remaining undefined. We continue in the simplest possible setting, one two-level system, so that two such systems can later be combined honestly.

\section{Eigenvalues Are Possible Measurement Results}

\begin{definition}
A nonzero vector \(\ket a\) is an eigenvector of \(M\) with eigenvalue \(\lambda_a\) if
\begin{equation}
  M\ket a=\lambda_a\ket a.
  \label{eq:l03-eigenvalue-equation}
\end{equation}
\end{definition}
The matrix may change the vector's length or phase, but it does not turn it into a different direction in state space.

For a Hermitian matrix the eigenvalue is real. Multiplying Equation~\eqref{eq:l03-eigenvalue-equation} on the left by \(\bra a\) gives
\begin{equation}
  \bra a M\ket a
  =\lambda_a\braket{a}{a}.
\end{equation}
The left side is real by Hermiticity, and the norm on the right is real and positive. Hence
\begin{equation}
  \lambda_a
  =
  \frac{\bra a M\ket a}{\braket{a}{a}}
  \in\mathbb R.
  \label{eq:l03-real-eigenvalue}
\end{equation}

Now comes the physical statement: when \(M\) represents an observable, its eigenvalues are the values that a measurement of \(M\) can return. The eigenvector belonging to \(\lambda_a\) is a state in which that value occurs with certainty.

The simplest example is a diagonal matrix:
\begin{equation}
  M=
  \begin{pmatrix}
    m_{11}&0\\
    0&m_{22}
  \end{pmatrix}.
\end{equation}

\begin{figure}[htbp]
\centering
\includegraphics[width=0.72\linewidth]{lecture_03_figure_03.png}
\caption{The abstract eigenvalue equation is placed above the first diagonal example.}
\lectureframe{00:44:04}
\end{figure}

Direct multiplication gives
\begin{align}
  M\begin{pmatrix}1\\0\end{pmatrix}
  &=m_{11}\begin{pmatrix}1\\0\end{pmatrix},
  \notag\\
  M\begin{pmatrix}0\\1\end{pmatrix}
  &=m_{22}\begin{pmatrix}0\\1\end{pmatrix}.
  \label{eq:l03-diagonal-eigenvectors}
\end{align}

\begin{figure}[htbp]
\centering
\includegraphics[width=0.72\linewidth]{lecture_03_figure_04.png}
\caption{The lower basis vector returns multiplied by the lower diagonal entry.}
\lectureframe{00:45:58}
\end{figure}

For a larger diagonal matrix, the same rule holds: its diagonal entries are eigenvalues and its standard one-hot columns are eigenvectors.

\section{The Three Spin Matrices}

For electron spin along the third axis, choose
\begin{equation}
  \sigma_3=
  \begin{pmatrix}
    1&0\\
    0&-1
  \end{pmatrix}.
\end{equation}
Its eigenstates are \(\ket u=(1,0)^T\) with eigenvalue \(+1\) and \(\ket d=(0,1)^T\) with eigenvalue \(-1\). Thus the two possible detector values are the two eigenvalues.

The first-axis spin matrix is
\begin{equation}
  \sigma_1=
  \begin{pmatrix}
    0&1\\
    1&0
  \end{pmatrix}.
\end{equation}
Its candidate eigenvectors are first written without normalization, since multiplying an eigenvector by any nonzero scalar leaves it an eigenvector. Physical states, however, must be normalized:
\begin{align}
  \ket{x+}
  &=\frac{1}{\sqrt2}\begin{pmatrix}1\\1\end{pmatrix},
  &
  \sigma_1\ket{x+}&=+\ket{x+},
  \notag\\
  \ket{x-}
  &=\frac{1}{\sqrt2}\begin{pmatrix}1\\-1\end{pmatrix},
  &
  \sigma_1\ket{x-}&=-\ket{x-}.
  \label{eq:l03-sigma-one-eigenvectors}
\end{align}

\begin{figure}[htbp]
\centering
\includegraphics[width=0.72\linewidth]{lecture_03_figure_05.png}
\caption{The off-diagonal matrix and equal-weight candidate appear before the product is completed.}
\lectureframe{00:50:29}
\end{figure}

\begin{classroomqa}
\lecturetimestamp{00:56:45}
\audiencequestion{Can the \(\sigma_1\) eigenvectors be seen geometrically rather than found by algebra?}
\lecturerresponse{In the real two-component picture, \(\sigma_1\) interchanges the basis directions \(E_1\) and \(E_2\). The diagonal direction \((1,1)^T\) is fixed, while the other diagonal \((1,-1)^T\) reverses. They are therefore eigenvectors with eigenvalues \(+1\) and \(-1\). This is a picture in component space, not yet a statement about the physical spin pointer.}
\end{classroomqa}

The remaining matrix is
\begin{equation}
  \sigma_2=
  \begin{pmatrix}
    0&-i\\
    i&0
  \end{pmatrix},
\end{equation}
with normalized eigenvectors
\begin{equation}
  \ket{y+}=\frac{1}{\sqrt2}\begin{pmatrix}1\\i\end{pmatrix},
  \qquad
  \ket{y-}=\frac{1}{\sqrt2}\begin{pmatrix}1\\-i\end{pmatrix}.
  \label{eq:l03-sigma-two-eigenvectors}
\end{equation}
Their eigenvalues are again \(+1\) and \(-1\).

All three matrices square to the identity:
\begin{equation}
  \sigma_1^2=\sigma_2^2=\sigma_3^2=I.
  \label{eq:l03-pauli-squares}
\end{equation}
If \(\sigma_i\ket a=\lambda\ket a\), then applying \(\sigma_i\) twice gives \(\lambda^2=1\), so \(\lambda=\pm1\). The identical spectra are the first sign that these matrices represent spin components along equivalent spatial axes.

\section{Distinct Outcomes Require Orthogonal States}

Before proving the next theorem, state its meaning. If one observable gives two different definite answers in two states, then a measurement of that observable can distinguish the states perfectly. Quantum mechanics expresses this perfect distinguishability as orthogonality.

\begin{theorem}
Let \(M\) be Hermitian, with
\begin{equation}
  M\ket A=\lambda_A\ket A,
  \qquad
  M\ket B=\lambda_B\ket B.
\end{equation}
If \(\lambda_A\ne\lambda_B\), then
\begin{equation}
  \braket{B}{A}=0.
\end{equation}
\end{theorem}

\begin{proof}
Multiplying the first eigenvalue equation by \(\bra B\) gives
\begin{equation}
  \bra B M\ket A=\lambda_A\braket{B}{A}.
\end{equation}
Hermiticity and the second eigenvalue equation give the same matrix element as
\begin{equation}
  \bra B M\ket A=\lambda_B\braket{B}{A},
\end{equation}
because \(\lambda_B\) is real. Subtraction yields
\begin{equation}
  (\lambda_A-\lambda_B)\braket{B}{A}=0.
\end{equation}
The eigenvalues are different, so the inner product must vanish.
\end{proof}

The \(\sigma_1\) eigenvectors provide an immediate check:
\begin{equation}
  \braket{x+}{x-}
  =
  \frac12
  \begin{pmatrix}1&1\end{pmatrix}
  \begin{pmatrix}1\\-1\end{pmatrix}
  =\frac12-\frac12=0.
\end{equation}
Conversely, for two orthogonal states there exists a measurement that distinguishes them. Orthogonality is therefore not merely a geometric word; it has an operational meaning.

\section{The Probability Postulate}

Suppose the system is prepared in a normalized state \(\ket B\). Let \(M\) be an observable, and let \(\ket A\) be a normalized eigenvector of \(M\) with eigenvalue \(\lambda_A\). Then the probability that measuring \(M\) returns \(\lambda_A\) is
\begin{equation}
  P(\lambda_A\mid B)
  =
  |\braket{A}{B}|^2
  =
  \braket{A}{B}\braket{B}{A}.
  \label{eq:l03-born-rule}
\end{equation}
The product is real and nonnegative. For normalized vectors it is at most one.

\begin{figure}[htbp]
\centering
\includegraphics[width=0.78\linewidth]{lecture_03_figure_06.png}
\caption{The probability rule is written as an overlap multiplied by its conjugate.}
\lectureframe{01:16:40}
\end{figure}

\begin{classroomqa}
\lecturetimestamp{01:19:19}
\audiencequestion{Is this the cosine squared of the angle between the two states?}
\lecturerresponse{For real unit vectors in ordinary Euclidean space, the inner product is the cosine of an angle. Here the inner product is generally complex, so there is no ordinary spatial angle whose cosine it is. The useful analogue is the magnitude of the overlap, and probability is its absolute square.}
\end{classroomqa}

\begin{classroomqa}
\lecturetimestamp{01:20:18}
\audiencequestion{What if the prepared state is another eigenvector of \(M\) with a different eigenvalue?}
\lecturerresponse{Then the two eigenvectors are orthogonal. Their overlap is zero, so the probability of obtaining \(\lambda_A\) is zero. Preparing a state with one definite value excludes a different definite value of the same observable.}
\end{classroomqa}

\begin{classroomqa}
\lecturetimestamp{01:21:23}
\audiencequestion{Can the same rule handle a preparation in which no outcome is certain?}
\lecturerresponse{Yes. That is the generic case: \(\ket B\) need not be an eigenstate of \(M\). Its overlaps with the normalized eigenvectors give the probabilities of the corresponding outcomes. We are beginning with a two-outcome example before treating more elaborate distributions.}
\end{classroomqa}

The lecture trains the rule on several preparations rather than hiding it inside one formula.

First prepare \(\sigma_3=+1\):
\begin{equation}
  \ket B=\ket{u}=\begin{pmatrix}1\\0\end{pmatrix}.
\end{equation}
Measure \(\sigma_1\). For its \(+1\) eigenvector,
\begin{equation}
  \braket{x+}{u}
  =
  \frac{1}{\sqrt2}
  \begin{pmatrix}1&1\end{pmatrix}
  \begin{pmatrix}1\\0\end{pmatrix}
  =\frac{1}{\sqrt2},
\end{equation}
so
\begin{equation}
  P(\sigma_1=+1\mid\sigma_3=+1)=\frac12.
\end{equation}
Replacing \(\ket{x+}\) by \(\ket{x-}\) gives the same probability. Reversing preparation and measurement also gives one half because the absolute square of the overlap is symmetric.

For a genuinely complex example, prepare \(\ket{x+}\) and measure \(\sigma_2\). Then
\begin{align}
  \braket{y+}{x+}
  &=
  \frac12
  \begin{pmatrix}1&-i\end{pmatrix}
  \begin{pmatrix}1\\1\end{pmatrix}
  =\frac{1-i}{2},
  \notag\\
  P(\sigma_2=+1\mid\sigma_1=+1)
  &=\left|\frac{1-i}{2}\right|^2
  =\frac14+\frac14
  =\frac12.
  \label{eq:l03-complex-probability-example}
\end{align}
Preparing spin along one axis and measuring along a perpendicular axis therefore gives equal probabilities for the two outcomes.

\section{Pointers, States, and Arbitrary Axes}

The preceding sentence can be misunderstood unless the vocabulary is repaired. A magnetic moment has a spatial direction, which we shall call a \emph{pointer}. A ket is a vector in the complex state space. Two spatial pointers may be perpendicular, while two state vectors may be orthogonal. These are different statements in different spaces.

\begin{classroomqa}
\lecturetimestamp{01:29:25}
\audiencequestion{If two physical spin directions are ninety degrees apart, should their state vectors be orthogonal?}
\lecturerresponse{No. Spin-up states along perpendicular spatial axes have overlap-squared \(1/2\), not zero. The state orthogonal to spin-up along one axis is spin-down along that same axis. We reserve \emph{perpendicular} for pointers in ordinary space and \emph{orthogonal} for perfectly distinguishable states.}
\end{classroomqa}

\begin{classroomqa}
\lecturetimestamp{01:31:25}
\audiencequestion{Could we use \emph{independent} instead of \emph{orthogonal}?}
\lecturerresponse{No. Linear independence only says that no vector is a scalar multiple of the others; it does not require zero inner products. Orthogonality is a stronger condition and here carries the physical meaning of perfect distinguishability.}
\end{classroomqa}

Let \(\hat n=(n_1,n_2,n_3)\) be a unit pointer in ordinary space:
\begin{equation}
  n_1^2+n_2^2+n_3^2=1.
\end{equation}
The operator for spin along that direction is the matrix-valued analogue of a dot product:
\begin{equation}
  \boldsymbol{\sigma}\cdot\hat n
  =
  n_1\sigma_1+n_2\sigma_2+n_3\sigma_3
  =
  \begin{pmatrix}
    n_3&n_1-in_2\\
    n_1+in_2&-n_3
  \end{pmatrix}.
  \label{eq:l03-spin-arbitrary-axis}
\end{equation}
It is Hermitian: the diagonal entries are real and the off-diagonal entries are conjugates.

\begin{figure}[htbp]
\centering
\includegraphics[width=0.84\linewidth]{lecture_03_figure_07.png}
\caption{The spatial unit pointer, the operator sum, and its explicit Hermitian matrix are developed together.}
\lectureframe{01:38:00}
\end{figure}

To find its possible values, square it before inserting matrix entries. Matrix order matters, so the cross terms occur as anticommutators:
\begin{align}
  (\boldsymbol{\sigma}\cdot\hat n)^2
  &=
  n_1^2\sigma_1^2+n_2^2\sigma_2^2+n_3^2\sigma_3^2
  \notag\\
  &\quad
  +n_1n_2(\sigma_1\sigma_2+\sigma_2\sigma_1)
  \notag\\
  &\quad
  +n_2n_3(\sigma_2\sigma_3+\sigma_3\sigma_2)
  \notag\\
  &\quad
  +n_1n_3(\sigma_1\sigma_3+\sigma_3\sigma_1).
  \label{eq:l03-axis-square-expanded}
\end{align}
Different Pauli matrices anticommute:
\begin{equation}
  \sigma_i\sigma_j+\sigma_j\sigma_i=0
  \qquad(i\ne j).
  \label{eq:l03-pauli-anticommutation}
\end{equation}
The lecture checks one case explicitly:
\begin{equation}
  \sigma_1\sigma_3
  =
  \begin{pmatrix}0&-1\\1&0\end{pmatrix},
  \qquad
  \sigma_3\sigma_1
  =
  \begin{pmatrix}0&1\\-1&0\end{pmatrix}
  =-\sigma_1\sigma_3.
  \label{eq:l03-order-check}
\end{equation}

\begin{figure}[htbp]
\centering
\includegraphics[width=0.82\linewidth]{lecture_03_figure_08.png}
\caption{Multiplying \(\sigma_1\) and \(\sigma_3\) in opposite orders produces matrices that differ by a minus sign.}
\lectureframe{01:44:50}
\end{figure}

The cross terms in Equation~\eqref{eq:l03-axis-square-expanded} vanish, while \(\sigma_i^2=I\). Hence
\begin{equation}
  (\boldsymbol{\sigma}\cdot\hat n)^2
  =(n_1^2+n_2^2+n_3^2)I
  =I.
\end{equation}
Every eigenvalue therefore satisfies \(\lambda^2=1\):
\begin{equation}
  \lambda=\pm1.
\end{equation}
Whichever spatial axis the apparatus selects, the detector still has only the two outcomes \(+1\) and \(-1\). The relation between nearby directions appears not as an intermediate measured value, but in the probabilities of those two fixed answers.
